\subsection{Joint multilayered LSIRM with item-only factor-mixture clustering}

\subsubsection{Index sets and basic notation}

Let $i=1,\dots,n$ index respondents, and let $\ell=1,\dots,L$ index data layers.
For layer $\ell$, let $\mathcal{J}_\ell$ denote the set of items belonging to that layer,
with cardinality $p_\ell = |\mathcal{J}_\ell|$.
Define the total number of items
\[
p = \sum_{\ell=1}^L p_\ell,
\]
and relabel the pooled item index as $j=1,\dots,p$.
Let $\ell(j)$ denote the layer to which item $j$ belongs.

For respondent $i$, let
\[
a_i \in \mathbb{R}^d
\]
denote the shared respondent latent position, and for item $j$, let
\[
b_j \in \mathbb{R}^d
\]
denote the item latent position in its corresponding layer.
In this paper we fix $d=2$.

Define the respondent--item latent distance
\[
d_{ij} = \|a_i - b_j\|, \qquad i=1,\dots,n,\; j=1,\dots,p.
\]

\subsubsection{LSIRM observation model}

The observation model for the multilayered LSIRM remains exactly the same as in the original
model specification.
That is, for each layer $\ell$, the conditional likelihood is written as
\[
p_\ell\!\left(Y_{ij}^{(\ell)} \mid \theta_{ij}^{(\ell)}\right),
\qquad
i=1,\dots,n,\; j\in\mathcal{J}_\ell,
\]
where the linear predictor depends on the shared respondent position $a_i$, the
item position $b_j$, and the layer-specific parameters.
For non-ordinal layers,
\[
\theta_{ij}^{(\ell)}
=
\alpha_i^{(\ell)} - \beta_j^{(\ell)} - \gamma^{(\ell)} d_{ij},
\]
and for ordinal layers the threshold-based predictor is defined as in the original multilayered LSIRM.

Hence the full LSIRM likelihood is
\[
p(Y \mid \Theta_{\mathrm{LSIRM}})
=
\prod_{\ell=1}^L \prod_{i=1}^n \prod_{j\in\mathcal{J}_\ell}
p_\ell\!\left(Y_{ij}^{(\ell)} \mid \theta_{ij}^{(\ell)}\right),
\]
where
\[
\Theta_{\mathrm{LSIRM}}
=
\left\{
a_i,\; b_j,\; \alpha_i^{(\ell)},\; \beta_j^{(\ell)},\; \gamma^{(\ell)},\;
\text{layer-specific auxiliary parameters}
\right\}.
\]

\subsubsection{Item relation profile induced by the current LSIRM state}

At each MCMC iteration, the current LSIRM configuration induces, for each item $j$,
a respondent-relation profile vector
\[
x_j
=
(x_{1j},\dots,x_{nj})^\top \in \mathbb{R}^n,
\]
where
\[
x_{ij} = \log d_{ij} = \log \|a_i-b_j\|.
\]

Thus the clustering layer is built directly on the current distance profile implied by the
LSIRM latent positions.

\subsubsection{Factor-mixture clustering layer}

Let $r \ll n$ denote the latent factor dimension, and let $K^\star$ be an overfitted upper bound
on the number of item clusters.

For each item $j$, introduce a latent factor score
\[
\eta_j \in \mathbb{R}^r,
\]
and a latent cluster label
\[
c_j \in \{1,\dots,K^\star\}.
\]

We model the respondent-relation profile of item $j$ by
\[
x_j \mid \delta,\Lambda,\eta_j,\Psi
\sim
N_n(\delta + \Lambda \eta_j,\Psi),
\]
where
\[
\delta = (\delta_1,\dots,\delta_n)^\top \in \mathbb{R}^n
\]
is a respondent-specific baseline vector,
\[
\Lambda \in \mathbb{R}^{n\times r}
\]
is the common factor loading matrix,
and
\[
\Psi = \mathrm{diag}(\psi_1,\dots,\psi_n)
\]
is a diagonal residual covariance matrix.

Conditional on its cluster label,
\[
\eta_j \mid c_j=l,\mu_l,\Sigma_l
\sim
N_r(\mu_l,\Sigma_l),
\qquad
l=1,\dots,K^\star.
\]

The cluster labels follow
\[
c_j \mid \rho \sim \mathrm{Categorical}(\rho),
\qquad
j=1,\dots,p,
\]
with mixture weights
\[
\rho = (\rho_1,\dots,\rho_{K^\star}) \sim \mathrm{Dirichlet}(e_0,\dots,e_0).
\]

\subsubsection{Prior specification for the clustering block}

For respondent-specific baselines,
\[
\delta_i \mid \sigma_\delta^2 \sim N(0,\sigma_\delta^2),
\qquad
i=1,\dots,n,
\]
\[
\sigma_\delta^2 \sim IG(a_\delta,b_\delta).
\]

Let $\lambda_i^\top$ denote the $i$th row of $\Lambda$.
Then
\[
\lambda_i \sim N_r(0,\tau_\lambda^2 I_r),
\qquad
i=1,\dots,n.
\]

For the diagonal residual variances,
\[
\psi_i \sim IG(a_\psi,b_\psi),
\qquad
i=1,\dots,n.
\]

For the Gaussian mixture in factor space,
\[
\mu_l \sim N_r(m_0,V_0),
\qquad
\Sigma_l \sim IW(\nu_0,S_0),
\qquad
l=1,\dots,K^\star.
\]

\subsubsection{Joint hierarchical model}

The full joint model is
\[
p(Y,\Theta_{\mathrm{LSIRM}},\Theta_{\mathrm{clust}})
=
p(Y \mid \Theta_{\mathrm{LSIRM}})
\,
p(X(\Theta_{\mathrm{LSIRM}})\mid \Theta_{\mathrm{clust}})
\,
p(\Theta_{\mathrm{LSIRM}})
\,
p(\Theta_{\mathrm{clust}}),
\]
where
\[
X(\Theta_{\mathrm{LSIRM}})
=
\{x_j(\Theta_{\mathrm{LSIRM}})\}_{j=1}^p,
\qquad
x_{ij}(\Theta_{\mathrm{LSIRM}})
=
\log \|a_i-b_j\|,
\]
and
\[
\Theta_{\mathrm{clust}}
=
\{\eta_j,c_j,\rho,\delta,\Lambda,\Psi,\mu_l,\Sigma_l,\sigma_\delta^2\}.
\]

Equivalently, the posterior kernel is
\[
p(\Theta_{\mathrm{LSIRM}},\Theta_{\mathrm{clust}} \mid Y)
\propto
p(Y \mid \Theta_{\mathrm{LSIRM}})
\left[
\prod_{j=1}^p
N_n\!\left(x_j(\Theta_{\mathrm{LSIRM}});\delta+\Lambda\eta_j,\Psi\right)
N_r\!\left(\eta_j;\mu_{c_j},\Sigma_{c_j}\right)
\rho_{c_j}
\right]
p(\Theta_{\mathrm{LSIRM}})
p(\rho)
p(\delta \mid \sigma_\delta^2)
p(\sigma_\delta^2)
p(\Lambda)
p(\Psi)
\prod_{l=1}^{K^\star} p(\mu_l)p(\Sigma_l).
\]

\subsubsection{Interpretation}

Under this model, item clustering is determined by the respondent-relation profiles
\[
x_j = (\log \|a_1-b_j\|,\dots,\log \|a_n-b_j\|)^\top,
\]
which are induced by the current LSIRM latent positions.
The matrix $\Lambda$ represents a shared low-dimensional basis for respondent-profile variation,
$\eta_j$ is the factor score of item $j$ in that basis, and the cluster label $c_j$
groups items with similar factor scores.

Importantly, there is no respondent clustering variable.
Thus respondent-side collapse does not force the item clustering to collapse into a trivial block model.

\subsubsection{Posterior updates for the clustering block}

Define
\[
\mathcal{J}_l = \{j : c_j=l\},
\qquad
n_l = |\mathcal{J}_l|.
\]

\paragraph{1. Update of $\rho$ (Gibbs)}
\[
\rho \mid c
\sim
\mathrm{Dirichlet}(e_0+n_1,\dots,e_0+n_{K^\star}).
\]

\paragraph{2. Update of $\eta_j$ (Gibbs)}
For each item $j$,
\[
\eta_j \mid \cdot
\sim
N_r(m_{\eta_j},V_{\eta_j}),
\]
where
\[
V_{\eta_j}
=
\left(
\Lambda^\top \Psi^{-1}\Lambda + \Sigma_{c_j}^{-1}
\right)^{-1},
\]
\[
m_{\eta_j}
=
V_{\eta_j}
\left(
\Lambda^\top \Psi^{-1}(x_j-\delta)
+
\Sigma_{c_j}^{-1}\mu_{c_j}
\right).
\]

\paragraph{3. Update of $c_j$ (categorical Gibbs)}
For each item $j$,
\[
P(c_j=l \mid \cdot)
\propto
\rho_l \,
\phi_r(\eta_j;\mu_l,\Sigma_l),
\qquad
l=1,\dots,K^\star,
\]
where $\phi_r(\cdot;\mu,\Sigma)$ denotes the $r$-dimensional Gaussian density.
Equivalently,
\[
\log \omega_{jl}
=
\log \rho_l
-
\frac{1}{2}\log |\Sigma_l|
-
\frac{1}{2}(\eta_j-\mu_l)^\top \Sigma_l^{-1}(\eta_j-\mu_l),
\]
and
\[
P(c_j=l \mid \cdot)
=
\frac{\exp(\log \omega_{jl})}
{\sum_{l'=1}^{K^\star}\exp(\log \omega_{jl'})}.
\]

\paragraph{4. Update of $\mu_l$ (Gibbs)}
For each cluster $l$,
\[
\mu_l \mid \cdot
\sim
N_r(m_{\mu_l},V_{\mu_l}),
\]
with
\[
V_{\mu_l}
=
\left(
V_0^{-1} + n_l \Sigma_l^{-1}
\right)^{-1},
\]
\[
m_{\mu_l}
=
V_{\mu_l}
\left(
V_0^{-1}m_0 + \Sigma_l^{-1}\sum_{j\in\mathcal{J}_l}\eta_j
\right).
\]

\paragraph{5. Update of $\Sigma_l$ (Gibbs)}
For each cluster $l$,
\[
\Sigma_l \mid \cdot
\sim
IW\left(
\nu_0+n_l,\;
S_0 + \sum_{j\in\mathcal{J}_l}(\eta_j-\mu_l)(\eta_j-\mu_l)^\top
\right).
\]

\paragraph{6. Update of respondent baseline $\delta_i$ (Gibbs)}
For each respondent $i$,
\[
\delta_i \mid \cdot
\sim
N(m_{\delta_i},V_{\delta_i}),
\]
where
\[
V_{\delta_i}
=
\left(
\frac{p}{\psi_i} + \frac{1}{\sigma_\delta^2}
\right)^{-1},
\]
\[
m_{\delta_i}
=
V_{\delta_i}
\left[
\frac{1}{\psi_i}
\sum_{j=1}^p
\bigl(x_{ij}-\lambda_i^\top \eta_j\bigr)
\right].
\]

\paragraph{7. Update of loading row $\lambda_i$ (Gibbs)}
For each respondent $i$,
\[
\lambda_i \mid \cdot
\sim
N_r(m_{\lambda_i},V_{\lambda_i}),
\]
where
\[
V_{\lambda_i}
=
\left(
\tau_\lambda^{-2}I_r
+
\psi_i^{-1}\sum_{j=1}^p \eta_j\eta_j^\top
\right)^{-1},
\]
\[
m_{\lambda_i}
=
V_{\lambda_i}
\left[
\psi_i^{-1}\sum_{j=1}^p \eta_j (x_{ij}-\delta_i)
\right].
\]

\paragraph{8. Update of $\psi_i$ (Gibbs)}
For each respondent $i$,
\[
\psi_i \mid \cdot
\sim
IG\left(
a_\psi + \frac{p}{2},
\;
b_\psi + \frac{1}{2}\sum_{j=1}^p
(x_{ij}-\delta_i-\lambda_i^\top \eta_j)^2
\right).
\]

\paragraph{9. Update of $\sigma_\delta^2$ (Gibbs)}
\[
\sigma_\delta^2 \mid \cdot
\sim
IG\left(
a_\delta + \frac{n}{2},
\;
b_\delta + \frac{1}{2}\sum_{i=1}^n \delta_i^2
\right).
\]

\subsubsection{Modified LSIRM updates under the joint model}

All layer-specific LSIRM updates for
\[
\alpha_i^{(\ell)},\; \beta_j^{(\ell)},\; \gamma^{(\ell)},
\]
and layer-specific auxiliary parameters remain the same as in the original multilayered LSIRM,
except that the updates of the latent positions $a_i$ and $b_j$ must now include the clustering-layer contribution.

\paragraph{1. Update of respondent position $a_i$ (MH)}
Suppose the proposal is
\[
a_i^\ast \sim q_a(\cdot \mid a_i).
\]
Then the acceptance probability is
\[
\alpha(a_i,a_i^\ast)
=
\min\left\{
1,\;
\frac{
p(Y \mid a_i^\ast,\Theta_{-a_i})
\, p(a_i^\ast)
\, \prod_{j=1}^p
N\!\left(\log\|a_i^\ast-b_j\|;\delta_i+\lambda_i^\top\eta_j,\psi_i\right)
\, q_a(a_i \mid a_i^\ast)
}{
p(Y \mid a_i,\Theta_{-a_i})
\, p(a_i)
\, \prod_{j=1}^p
N\!\left(\log\|a_i-b_j\|;\delta_i+\lambda_i^\top\eta_j,\psi_i\right)
\, q_a(a_i^\ast \mid a_i)
}
\right\}.
\]

\paragraph{2. Update of item position $b_j$ (MH)}
Suppose the proposal is
\[
b_j^\ast \sim q_b(\cdot \mid b_j).
\]
Then the acceptance probability is
\[
\alpha(b_j,b_j^\ast)
=
\min\left\{
1,\;
\frac{
p(Y_{\cdot j} \mid b_j^\ast,\Theta_{-b_j})
\, p(b_j^\ast)
\, \prod_{i=1}^n
N\!\left(\log\|a_i-b_j^\ast\|;\delta_i+\lambda_i^\top\eta_j,\psi_i\right)
\, q_b(b_j \mid b_j^\ast)
}{
p(Y_{\cdot j} \mid b_j,\Theta_{-b_j})
\, p(b_j)
\, \prod_{i=1}^n
N\!\left(\log\|a_i-b_j\|;\delta_i+\lambda_i^\top\eta_j,\psi_i\right)
\, q_b(b_j^\ast \mid b_j)
}
\right\}.
\]
Here $Y_{\cdot j}$ denotes the subset of observed responses associated with item $j$
in its corresponding layer.

\subsubsection{One full MCMC sweep}

One full MCMC iteration proceeds as follows:
\[
\text{LSIRM non-position block}
\;\rightarrow\;
\{a_i\}_{i=1}^n
\;\rightarrow\;
\{b_j\}_{j=1}^p
\;\rightarrow\;
\rho
\;\rightarrow\;
\{\eta_j\}_{j=1}^p
\;\rightarrow\;
\{c_j\}_{j=1}^p
\;\rightarrow\;
\{\mu_l,\Sigma_l\}_{l=1}^{K^\star}
\;\rightarrow\;
\{\delta_i,\lambda_i,\psi_i\}_{i=1}^n
\;\rightarrow\;
\sigma_\delta^2.
\]

\subsubsection{Pseudo-code}

\begin{algorithm}[H]
\caption{Joint multilayered LSIRM + item-only factor-mixture clustering}
\begin{algorithmic}[1]
\Require Observed multilayer item response data $\{Y_{ij}^{(\ell)}\}$, factor dimension $r$,
overfitted number of clusters $K^\star$, total iterations $T$, burn-in $B$, thinning $H$.
\State Initialize LSIRM parameters
$\Theta_{\mathrm{LSIRM}}^{(0)}$
and clustering parameters
$\Theta_{\mathrm{clust}}^{(0)}$.
\For{$t=1,\dots,T$}
    \State Update all layer-specific LSIRM parameters except positions
    using the original multilayered LSIRM full conditionals or MH steps.
    \For{$i=1,\dots,n$}
        \State Update respondent position $a_i$ by MH using the original LSIRM likelihood
        multiplied by the clustering-layer Gaussian contribution.
    \EndFor
    \For{$j=1,\dots,p$}
        \State Update item position $b_j$ by MH using the original item-specific LSIRM likelihood
        multiplied by the clustering-layer Gaussian contribution.
    \EndFor
    \State Compute current distances
    \[
    d_{ij}^{(t)} = \|a_i^{(t)}-b_j^{(t)}\|,
    \qquad
    x_{ij}^{(t)} = \log d_{ij}^{(t)}.
    \]
    \State Update $\rho$ from its Dirichlet full conditional.
    \For{$j=1,\dots,p$}
        \State Update factor score $\eta_j$ from its Gaussian full conditional.
    \EndFor
    \For{$j=1,\dots,p$}
        \State Update cluster label $c_j$ from its categorical full conditional.
    \EndFor
    \For{$l=1,\dots,K^\star$}
        \State Update $\mu_l$ from its Gaussian full conditional.
        \State Update $\Sigma_l$ from its inverse-Wishart full conditional.
    \EndFor
    \For{$i=1,\dots,n$}
        \State Update respondent baseline $\delta_i$ from its Gaussian full conditional.
        \State Update loading row $\lambda_i$ from its Gaussian full conditional.
        \State Update residual variance $\psi_i$ from its inverse-gamma full conditional.
    \EndFor
    \State Update $\sigma_\delta^2$ from its inverse-gamma full conditional.
    \If{$t>B$ and $(t-B)$ is divisible by $H$}
        \State Store the current draw of
        $\Theta_{\mathrm{LSIRM}}^{(t)}$ and $\Theta_{\mathrm{clust}}^{(t)}$.
    \EndIf
\EndFor
\end{algorithmic}
\end{algorithm}

\subsubsection{Posterior summary}

The occupied number of item clusters at iteration $t$ is
\[
K_+^{(t)} = \sum_{l=1}^{K^\star} I(n_l^{(t)} > 0).
\]
The posterior co-clustering probability between items $j$ and $j'$ is
\[
\widehat P_{jj'}
=
\frac{1}{T_{\mathrm{save}}}
\sum_{t=1}^{T_{\mathrm{save}}}
I\!\left(c_j^{(t)} = c_{j'}^{(t)}\right).
\]
A representative final partition may be obtained from the posterior similarity matrix
$\widehat P = (\widehat P_{jj'})$ after standard label-switching post-processing.